\section{Research Lineage: V6 through V6.7}
\label{sec:archive-lineage}

The earlier versions investigated several objects that are related to physical reconstruction but have distinct
targets: reproducible editing, occupancy authority, conditional state selection, trajectory reliability,
and local surface repair. This section retains their strongest results and terminal comparisons.
The version closeouts and technical reports under \path{docs/autoresearch/worldsim_v6*} are the evidence sources.

\subsection{V6: Typed World State and Exact Reuse}

V6 established SceneIR ownership, deterministic world editing, provenance, lifecycle handling, and factorized
verification. It separated exact reconstruction of a frozen output from the semantic or physical validity of an edit.
Its final selector family studied when perception could be reused without changing the full pipeline's outputs.

The R139 orthogonal confirmation reduced perception calls by 17.95\%, with zero false reuse and
156 exact output hashes. R140 included shared sensor rendering in both full and selective costs:
end-to-end reductions were 13.53\% for StreetGS, 11.14\% for AD-GS development, and 1.66\% for
AD-GS confirmation, averaging 8.78\% across conditions with zero reconstruction errors.
The difference between call-count and end-to-end savings shows the effect of shared rendering cost.

The threshold-transfer alternatives remained negative: R134 had an AD-GS false negative and R136 a
held-out conservative false positive. The identity-only guard was retained instead of a universal
pixel-threshold claim. R140's first two closeouts failed through Python/JSON boolean handling;
the third reproduced their scientific certificate and completed the engineering record.
R141 was not run. This lineage contributes state ownership and exact operational reproducibility to
the current method, while its selector is outside the EAS inference path.

\subsection{V6.1: Occupancy as a Physical Evidence Source}

The 28-case minimum experiment separated method-visible occupancy from hidden observed-FREE evaluation.
Oracle occupancy accepted 10/28 cases with zero false-safe cases and 39.83\% mask yield.
This established useful proposal support under the oracle.

Hunyuan's generic generated surfaces accepted 0/6 cases in each tested arm and conflicted with measured free space.
Two independent predicted-occupancy backends, GaussianWorld and IR-WM, each recovered the oracle's 10/28 support
but produced 10/10 false-safe accepted cases.
The result concerns the frozen uncalibrated argmax-occupancy interfaces, not occupancy learning in general.
ME-4 was unexecuted after the terminal comparison.
The retained lesson for EAS is to keep unknown and contradictory observations explicit in the supervised target.

\subsection{V6.2--V6.3: Constraints and Surface Risk}

V6.2 introduced a constraint-aware evidential state and a single evidence-dropout recovery.
With $p=.5$ feature dropout and teacher consistency, source-valid UNKNOWN fell from 82.74\% to 63.85\%.
The four accepted legacy cases nevertheless remained the same four false-safe cases.
The recovered projection had zero hard violations over 939,206 queries and full safe-OCC retention.
Thus satisfying known constraints and reducing abstention did not resolve hidden free-space errors.

V6.3 moved to native surface and tail-risk supervision.
Early objectives collapsed to infeasible states; a recovered checkpoint made the common comparison executable.
On two development scenes, the B3 surface model increased pooled tail CVaR from .491496 to .608174
while reducing emitted occupied points from 2,298,450 to 1,047,186.
Both scene-level comparisons failed. Zero known hard violations and nonzero useful output coexisted
with worse hidden-surface risk. Surface-Max, Surface-CVaR, and downstream authority arms remained unexecuted.
These results motivate direct supervision of both the surface and the measured return event.

\subsection{V6.4: Conditional Compilation at Fixed Opportunity}

Native uncertainty features became useful after a full 273-dimensional risk head replaced the weaker
PCA-16 representation. At nominal .4 coverage, independent calibration had zero failures and
confirmation had one failure in 96 cases at .39994 realized coverage.
A stratum-conditional map raised mean coverage to approximately .475 with zero failures on its
separate 96-case confirmation cohort.

The strongest route-aware result uses an untouched, fixed-opportunity 96-case denominator:
M1 preserved M0 mean coverage .474970 while reducing worst-ten route-conflict CVaR
from .020726 to .010821 and pooled conflict density from .004945 to .002001.
M1 was lower/equal/higher in 18/78/0 paired cases.
Earlier selected-count ratios were kept separate because they changed the evaluated opportunity set.

The collision-critic extension failed. Initial unsafe-action recalls were .02174/0/.01087 for
Real-only/naive/verified augmentation on 184 unsafe actions.
Independent threshold recovery achieved .62044 recall for the verified arm and still produced two
selected-policy false-safes; AUROC shifted from .71165 to .56274.
The relative route-conflict result therefore remains an earlier positive experiment with its own target.
It is not evidence of a learned collision controller.

\subsection{V6.5: Visited-State Reliability}

V6.5 changed the query to the reliability of world states visited by a given Ego trajectory over two seconds.
Frozen Qmean transferred on fresh cohorts: P2V obtained Spearman .63396, unsafe AUROC .99415,
and 49.25\% lower selected cost.
A two-parameter monotone calibration map reduced MSE by 92.80\% and five-bin calibration error by
88.31\% without changing ranking. P10V's fixed-action cohort retained Spearman .74023,
AUROC .85878, and 33.26\% lower selected cost.

The combined P10X confirmation passed five of six criteria, but selected cost decreased only 16.38\%
against the frozen 25\% requirement. This was a terminal rejection of the direct action-selection benefit.
The map, Actor false-safe, learned-admission, temporal, and smooth-tail branches remain recorded as
unsuccessful alternatives. The retained positive is a trajectory-conditioned reliability evaluator.

\subsection{V6.6: Actor Preservation and Local Repair}

V6.6 separated Actor existence from local geometry validity.
The frozen local geometry head reached AUROC/AUPRC .761644/.767165 on an independent consumed
legacy cohort covering six scenes. A runtime package retained 127 Actors, 581 states, and 1.62 million
primitives. Fixed-budget triage reduced untreated conflict exposure by 68.40\%.

Natural surface repair remained a trade-off.
Exact sensor support reduced conflicts by 84.77\% but retained only 39.57\% of clean geometry.
One-voxel support retained 61.95\% of clean geometry but reduced conflicts by only 41.79\%.
The single allowed recovery failed the joint criterion.
Separately, a corrected bounded responder passed six synthetic lead-brake cases, reducing collision
steps from 306 to zero with jerk at most six. This is a synthetic capability result.
Natural repair and downstream learned control were not established by that demonstration.

\subsection{V6.7: Uncertainty, Task Geometry, and Reliability Surfaces}

V6.7 developed the earlier paper's reliability formulation through a long sequence of query objects.
The central successful factorization predicts Actor residual uncertainty from source state and supplies
candidate-trajectory clearance analytically.
Endpoint prediction error itself is independent of which candidate trajectory queries it; task geometry
determines the decision boundary and which errors matter.

\paragraph{From endpoint summaries to directional uncertainty.}
P81 reduced fixed-50\% unreliable events from 57 to 26 on ten fresh scenes.
However, the P84--P94 summary, Deep Sets, attention, ordinal, continuous, Gaussian, BCE, and ensemble
variants failed to stably improve trajectory-level visited-state objectives.
Task-conditioned occupancy-flip classifiers also reversed across cohorts.
The uncertainty-tube factorization reduced fixed-50\% flips from 35 to five on the fresh P108 cohort.
Directional Gaussian projection strengthened ranking, while P111 exposed deterministic clearance as a strong
comparator. Subsequent work therefore compared learned uncertainty directly with geometry.

P147 confirmed a multi-horizon continuous-cost result on ten fresh scenes:
all five horizons (.8/1.5/2.5/3.0/3.5 seconds) had higher rank correlation and lower selected cost.
Macro Spearman increment was .17419 and selected-cost difference was $-.01777$.
This is distinct from the earlier binary event and from an absolute probability estimate.

\paragraph{Conditional density and joint horizons.}
P182 modeled $\log(1+C)$ with a five-component conditional Gaussian mixture.
P183's separate ten-scene confirmation reduced integrated Brier by 28.48\% and mean absolute
reliability error by 69.38\% over P173.
P199 then coupled aligned horizon marginals through input-conditioned Gaussian dependence.
On 1,846 trajectories in the ten-scene P201 cohort, Brier improved by 17.52\% and calibration error
by 53.37\% relative to independent marginals.
The P203 beta map received same-read prospective secondary support, improving P199 Brier by 3.70\%
and calibration error by 49.12\%.

Alternative sampling, density, and dependence branches are retained as comparisons:
bootstrap/DRO, condition smoothing, heavy-tailed density, spline flows, direct Brier optimization,
scene-balanced sampling, mixtures, constant or low-rank copulas, Student-t copulas, direct maximum-cost
density, prefix-survival density, cumulative exposure, and learned selection.
Several improved development calibration or one cohort while regressing another.
Their results did not replace the independently supported density/dependence stages.

\paragraph{Monotone runtime distillation.}
P227 distilled the frozen density/dependence teacher.
P228's fresh 1,720-trajectory cohort had teacher MAE .007443, a .316\% Brier improvement,
and a .000020 absolute calibration increase.
Its measured batched student stage took .002905 seconds versus .111554 seconds for the teacher joint stage
(38.4 times at that stage, excluding shared upstream work).
P234's same-read secondary surface had MAE .007101 and zero budget/horizon monotonicity violations.
Input-reduction and global continuous-family variants missed their fidelity criteria.
Fixed-knot interpolation retained the discrete surface and its structural monotonicity.

\paragraph{Terminal selective-authority family.}
P346 retained a reused-development result: q90 mean coverage .328415 with maximum unsafe selected-set
rate .090909 on P201 task conditions. Its source-held-out-horizon rate was .267108.
The executable lattice's held-out horizon is 2.5 seconds; V7-F16 corrects the older prose label of 3.0 seconds.
P358--P361 then trained all four fixed-grid horizons and tested shared calibration, horizon calibration,
one-sided BCE, and pairwise ranking. Their P201 coverage/risk pairs were
.306284/.112360, .293169/.105263, .342896/.107527, and .343169/.127660.
None simultaneously met coverage at least .30 and risk at most .10.
P346 therefore remains a reused-development positive; the terminal grid family remains negative.

\subsection{Transition to V7 and the Current Surface Object}

V7 combined Actor-preserving physical compilation, separate repair/hazard inputs, and the inherited reliability
surface. The non-learned AV2 compiler improved depth and Chamfer while preserving Actor state.
Selective repair transferred on some cohorts, but fresh ranking reversal, sensor-opportunity shortcuts,
and confidently wrong interval-stable decisions exposed its limitations.
The fixed-lattice downstream audit checked how repair and trajectory reliability compose as two independent
decisions; it was an open-loop interface study.

V7.1 moved the primary prediction object to supervised canonical surface geometry and return evidence.
The current model inherits explicit state ownership and disciplined denominator accounting.
The old selector, task-reliability density, and synthetic responder remain useful historical results in
this archive, with their full V7 tables and proof record below. They are not hidden components of EAS.
